\chapter{Hypothermia}
\index{Hypothermia}

\section*{Definition and Overview}
Hypothermia is a medical emergency in which the body's core temperature falls below 35 degrees Celsius, and the body loses heat faster than it can produce it. It occurs with prolonged exposure to cold, especially with wet clothing, wind, or immersion in cold water. Mild hypothermia causes shivering and confusion, while severe hypothermia leads to loss of consciousness and can be fatal. Older adults, babies, people with certain medical conditions, and those outdoors or homeless are at highest risk. Prompt rewarming and emergency care are essential.

\section*{Epidemiology}
Hypothermia is uncommon in many countries compared with colder countries, but it occurs during winter, after cold-water immersion, in outdoor activities, and in vulnerable people such as the elderly living in cold homes. It contributes to deaths in cold weather and after accidents, and it can complicate alcohol intoxication and drug use. Older people with reduced temperature regulation are at particular risk, and hypothermia in hospitalised and postoperative patients is also recognised.

\section*{Aetiology and Pathophysiology}
\begin{itemize}
    \item \textbf{Cold exposure:} cold air, wind, rain, and cold water
    \item \textbf{Wet clothing:} water conducts heat away rapidly
    \item \textbf{Immersion:} cold-water immersion causes rapid heat loss
    \item \textbf{Reduced heat production:} exhaustion, malnutrition, and alcohol
    \item \textbf{Impaired regulation:} older age, babies, and medical conditions
    \item \textbf{Medications and drugs:} alcohol and sedatives impair temperature control
    \item \textbf{Pathophysiology:} as core temperature falls, shivering increases heat production, then fails; metabolism slows, and organ function deteriorates
    \item \textbf{Cardiac risk:} hypothermia predisposes to dangerous heart rhythms
\end{itemize}

\section*{Clinical Features}
\begin{itemize}
    \item Shivering, which stops in severe hypothermia
    \item Cold, pale skin and slurred speech
    \item Confusion, drowsiness, and poor coordination
    \item Slow breathing and slow pulse
    \item Weakness and fatigue
    \item Loss of consciousness in severe cases
    \item In babies: cold skin, lethargy, and poor feeding
    \item Paradoxical undressing in severe cases
\end{itemize}

\section*{Differential Diagnosis}
\begin{itemize}
    \item Stroke and other neurological conditions can mimic hypothermia
    \item Hypoglycaemia causes confusion and drowsiness
    \item Sepsis and other severe illness
    \item Drug and alcohol intoxication
    \item Core temperature measurement confirms the diagnosis
\end{itemize}

\section*{When to Seek Medical Care}
Anyone with suspected hypothermia needs urgent medical assessment. Move the person to warmth, remove wet clothing, and call for help.

\textbf{Contact your local emergency services for:}
\begin{itemize}
    \item Confusion, drowsiness, or loss of consciousness in the cold
    \item Shivering that stops with persistent coldness
    \item A body temperature below 35 degrees Celsius
    \item Slow breathing or slow pulse
    \item A baby who is cold and lethargic
    \item Cold-water immersion or prolonged cold exposure
\end{itemize}

\section*{Diagnosis and Investigations}
\begin{itemize}
    \item \textbf{Core temperature measurement:} a low-reading thermometer confirms the diagnosis
    \item \textbf{ECG:} for the heart rhythm changes of hypothermia
    \item \textbf{Blood tests:} blood sugar, electrolytes, and organ function
    \item \textbf{Assessment of the cause:} exposure, alcohol, illness, and medications
    \item \textbf{Monitoring:} temperature, heart rhythm, and consciousness during rewarming
\end{itemize}

\section*{Management}
\begin{itemize}
    \item \textbf{Mild hypothermia:} move to warmth, remove wet clothes, and use blankets and warm drinks
    \item \textbf{Passive rewarming:} insulation for mild cases
    \item \textbf{Active rewarming:} warm fluids, heating blankets, and warm air in hospital
    \item \textbf{Severe hypothermia:} intensive care with careful rewarming
    \item \textbf{Cardiac care:} monitoring and treatment of arrhythmias
    \item \textbf{Handle gently:} rough movement can trigger dangerous rhythms
    \item \textbf{Treat the cause:} alcohol, hypoglycaemia, and underlying illness
    \item \textbf{"Not dead until warm and dead":} resuscitation continues until rewarmed
\end{itemize}

\section*{Complications}
\begin{itemize}
    \item Cardiac arrest and dangerous arrhythmias
    \item Brain injury from prolonged cold
    \item Frostbite of the extremities
    \item Pneumonia and other infections
    \item Kidney injury and blood clotting problems
    \item Falls and injuries during confusion
    \item Death in severe untreated cases
\end{itemize}

\section*{Prognosis and Outcomes}
Mild hypothermia resolves fully with rewarming. Even severe hypothermia has a surprisingly good outlook with appropriate care, because the cold protects organs during resuscitation; people have recovered after prolonged cardiac arrest in hypothermia. The outcome depends on the depth and duration of cold exposure and the speed of treatment. With prevention and prompt care, the risks of hypothermia are largely avoidable.

\section*{Prevention and Self-Care}
\begin{itemize}
    \item Dress warmly in layers and wear waterproof clothing in wet or windy conditions
    \item Avoid prolonged cold exposure and take regular warm breaks
    \item Never drink alcohol in cold conditions
    \item Check on older relatives and neighbours during cold weather
    \item Keep babies and older people warm
    \item Plan outdoor activities with cold weather in mind
    \item Carry warm clothing and emergency supplies on trips
    \item Learn first aid for hypothermia
\end{itemize}

\section*{Sources and Further Reading}
The following reputable sources provide further information for healthcare students and patients seeking reliable, current advice about hypothermia.

\begin{itemize}
    \item Healthdirect. \textit{Hypothermia: first aid.} https://www.healthdirect.gov.au/hypothermia
    \item St John Ambulance Australia. \textit{First aid for hypothermia.} https://stjohn.org.au/
    \item Better Health Channel. \textit{Hypothermia.} https://www.betterhealth.vic.gov.au/
    \item NHS. \textit{Hypothermia.} https://www.nhs.uk/conditions/hypothermia/
    \item Mayo Clinic. \textit{Hypothermia: first aid.} https://www.mayoclinic.org/
    \item MedlinePlus. \textit{Hypothermia.} https://medlineplus.gov/hypothermia.html
\end{itemize}
